\documentclass[12pt]{article}

\usepackage{sbc-template}
\usepackage{graphicx,url}
\usepackage[utf8]{inputenc}
\usepackage[english]{babel}

\sloppy

\title{Relatório Prático: Kernel em RISC-V\\
\large Alocação de Memória Física, Processos e Paginação Sv32}

\author{João Pedro Sant' Ana de Carvalho}

\address{Faculdade de Computação -- Universidade Federal de Mato Grosso do Sul (UFMS)\\
  Campo Grande -- MS -- Brasil
  \email{joao.sant@ufms.br}
}

\begin{document}

\maketitle

\begin{abstract}
This report documents the continuation of the minimal 32-bit RISC-V operating system kernel, focusing on physical memory management, multitasking support, and virtual memory paging. The completed tasks include dynamic page allocation via a bump allocator with linker-defined memory bounds (Task 6), process infrastructure featuring process control blocks, private stacks, context switching in assembly, cooperative scheduling, and trap stack protection via the sscratch register (Task 7), as well as two-level Sv32 page table management, kernel identity mapping, and MMU address translation via the satp register (Task 8). Tasks 9 (user mode isolation) and 10 (system calls) were not implemented in this version. The implementation was validated on the QEMU emulator.
\end{abstract}

\begin{resumo}
Este relatório descreve a evolução de um kernel mínimo para a arquitetura RISC-V de 32 bits, com foco no gerenciamento de memória física, no suporte a múltiplos fluxos de execução e na introdução da memória virtual. Foram implementados o alocador dinâmico de páginas com delimitação de memória livre via linker script (Trabalho 06), a infraestrutura de processos com blocos de controle (PCB), pilhas privativas, troca de contexto em assembly, escalonamento cooperativo e proteção da pilha de interrupção com o registrador sscratch (Trabalho 07), além das tabelas de páginas em dois níveis no modelo Sv32, mapeamento identidade da região do kernel e chaveamento da MMU via registrador satp com barreiras sfence.vma (Trabalho 08). Fica registrado que os Trabalhos 09 (modo usuário) e 10 (chamadas de sistema) não foram realizados nesta entrega. As implementações foram validadas no emulador QEMU.
\end{resumo}

\section{Introdução}

Após o estabelecimento do boot em modo \textit{bare-metal}, da saída serial via chamadas OpenSBI e do suporte a exceções da CPU, o desenvolvimento do sistema avançou para os pilares essenciais de gerenciamento de memória e concorrência em um sistema operacional.

Este relatório abrange as atividades concluídas e testadas correspondentes ao escopo da segunda entrega prática da disciplina:
\begin{itemize}
    \item \textbf{Trabalho 06}: Alocação de memória física alinhada a páginas de 4~KB;
    \item \textbf{Trabalho 07}: Estrutura de processos, troca de contexto e escalonamento cooperativo;
    \item \textbf{Trabalho 08}: Tabelas de páginas Sv32 e gerenciamento de memória virtual.
\end{itemize}

As tarefas subsequentes previstas para o ciclo da Entrega 2 (\textbf{Trabalhos 09 e 10}), que englobam a transição para o nível de privilégio de usuário e a interface de chamadas de sistema (\texttt{ecall}), \textbf{não foram realizadas} nesta etapa e permanecem como trabalho futuro.

\section{Atividades Realizadas}

\subsection{Trabalho 06: Alocação de Memória Física}

Para permitir a criação dinâmica de estruturas pelo kernel sem recorrer à alocação estática arbitrária, foi implementado um alocador de páginas físicas baseado na estratégia \textit{Bump Allocator} (ou alocador linear).

A delimitação da memória disponível foi realizada diretamente no linker script (\texttt{kernel.ld}), logo após o espaço reservado para a pilha do kernel (\texttt{\_\_stack\_top}). A área foi alinhada a uma fronteira de 4096 bytes (4~KB) e estendida em 64~MB, demarcada pelos símbolos \texttt{\_\_free\_ram} e \texttt{\_\_free\_ram\_end}:

\begin{verbatim}
    . = ALIGN(4096);
    __free_ram = .;
    . += 64 * 1024 * 1024;
    __free_ram_end = .;
\end{verbatim}

No arquivo \texttt{common.h}, foi definida a constante fundamental \texttt{PAGE\_SIZE} como 4096. A função \texttt{alloc\_pages} no arquivo \texttt{kernel.c} gerencia a alocação de $n$ páginas consecutivas. A variável estática \texttt{next\_paddr} mantém o endereço da próxima região livre. A inicialização é resolvida dinamicamente no primeiro ciclo para garantir estrita conformidade com expressões constantes em C:

\begin{verbatim}
paddr_t alloc_pages(uint32_t n) {
  static paddr_t next_paddr = 0;
  if (!next_paddr)
    next_paddr = (paddr_t)__free_ram;

  paddr_t paddr = next_paddr;
  next_paddr += n * PAGE_SIZE;

  if (next_paddr > (paddr_t)__free_ram_end)
    PANIC("out of memory");

  memset((void *)paddr, 0, n * PAGE_SIZE);
  return paddr;
}
\end{verbatim}

A função assegura três invariantes:
\begin{enumerate}
    \item Os endereços retornados são múltiplos exatos de 4~KB;
    \item Em caso de esgotamento do limite superior (\texttt{\_\_free\_ram\_end}), dispara-se \texttt{PANIC("out of memory")};
    \item A memória retornada é integralmente limpa com zeros via \texttt{memset}, evitando a propagação de lixo de memória.
\end{enumerate}

\subsection{Trabalho 07: Processos e Troca de Contexto}

O sétimo trabalho viabilizou a execução concorrente de tarefas no kernel. No arquivo \texttt{kernel.h}, foi definido o Bloco de Controle de Processo (\texttt{struct process}) e a capacidade máxima de processos no sistema (\texttt{PROCS\_MAX = 8}):

\begin{verbatim}
#define PROCS_MAX 8
#define PROC_UNUSED   0
#define PROC_RUNNABLE 1

struct process {
  int pid;
  int state;
  vaddr_t sp;
  uint32_t *page_table;
  uint8_t stack[8192];
};
\end{verbatim}

Cada processo possui sua própria pilha de execução com 8~KB de extensão, eliminando a dependência de pilha única global.

A rotina de troca de contexto (\texttt{switch\_context}) foi implementada em assembly com o atributo \texttt{naked}. Seguindo a convenção de chamada da ABI RISC-V, apenas os 13 registradores \textit{callee-saved} (\texttt{ra} e \texttt{s0} a \texttt{s11}) necessitam de salvamento na pilha da tarefa corrente, uma vez que os registradores temporários (\textit{caller-saved}) já são geridos pelo compilador no momento da invocação:

\begin{verbatim}
__attribute__((naked)) void switch_context(uint32_t *prev_sp,
                                           uint32_t *next_sp) {
  __asm__ __volatile__(
      "addi sp, sp, -13 * 4\n"
      "sw ra,  0  * 4(sp)\n"
      "sw s0,  1  * 4(sp)\n"
      "sw s1,  2  * 4(sp)\n"
      "sw s2,  3  * 4(sp)\n"
      "sw s3,  4  * 4(sp)\n"
      "sw s4,  5  * 4(sp)\n"
      "sw s5,  6  * 4(sp)\n"
      "sw s6,  7  * 4(sp)\n"
      "sw s7,  8  * 4(sp)\n"
      "sw s8,  9  * 4(sp)\n"
      "sw s9,  10 * 4(sp)\n"
      "sw s10, 11 * 4(sp)\n"
      "sw s11, 12 * 4(sp)\n"
      "sw sp, (a0)\n"
      "lw sp, (a1)\n"
      "lw ra,  0  * 4(sp)\n"
      "lw s0,  1  * 4(sp)\n"
      "lw s1,  2  * 4(sp)\n"
      "lw s2,  3  * 4(sp)\n"
      "lw s3,  4  * 4(sp)\n"
      "lw s4,  5  * 4(sp)\n"
      "lw s5,  6  * 4(sp)\n"
      "lw s6,  7  * 4(sp)\n"
      "lw s7,  8  * 4(sp)\n"
      "lw s8,  9  * 4(sp)\n"
      "lw s9,  10 * 4(sp)\n"
      "lw s10, 11 * 4(sp)\n"
      "lw s11, 12 * 4(sp)\n"
      "addi sp, sp, 13 * 4\n"
      "ret\n");
}
\end{verbatim}

A função \texttt{create\_process} aloca um descritor livre na tabela estática \texttt{procs}, monta a estrutura inicial da pilha empilhando zeros para os registradores salvos e o endereço de entrada no registrador \texttt{ra}, ajustando o estado para \texttt{PROC\_RUNNABLE}.

O escalonador cooperativo (\texttt{yield}) realiza a busca circular em rodízio pelo próximo processo executável com \texttt{pid > 0}. Caso nenhum outro processo esteja pronto, o \texttt{idle\_proc} (PID 0) é selecionado. Durante a transição, a base da pilha do próximo processo é gravada no registrador CSR \texttt{sscratch}, permitindo que o tratador \texttt{kernel\_entry} recupere de forma segura a pilha privativa do processo através da instrução \texttt{csrrw sp, sscratch, sp}.

\subsection{Trabalho 08: Tabelas de Páginas e Memória Virtual Sv32}

O oitavo trabalho introduziu o suporte ao mecanismo de tradução e isolamento de memória virtual Sv32 do RISC-V. No modelo Sv32, os endereços virtuais de 32 bits são decompostos em dois níveis hierárquicos: o índice de primeiro nível ($VPN[1]$, 10 bits), o índice de segundo nível ($VPN[0]$, 10 bits) e o deslocamento dentro da página ($Offset$, 12 bits).

No arquivo \texttt{kernel.h}, foram estabelecidas as definições de controle do registrador \texttt{satp} e das entradas da tabela de páginas (PTE):
\begin{verbatim}
#define SATP_SV32 (1u << 31)
#define PAGE_V    (1 << 0) // Entrada válida
#define PAGE_R    (1 << 1) // Permissão de leitura
#define PAGE_W    (1 << 2) // Permissão de escrita
#define PAGE_X    (1 << 3) // Permissão de execução
#define PAGE_U    (1 << 4) // Acesso em modo usuário
\end{verbatim}

A estrutura de processo (\texttt{struct process}) foi atualizada para conter o ponteiro da tabela raiz de primeiro nível (\texttt{uint32\_t *page\_table}).

No linker script (\texttt{kernel.ld}), exportou-se o símbolo \texttt{\_\_kernel\_base = .;} logo após a atribuição do endereço base em \texttt{0x80200000}, estabelecendo a fronteira inferior da memória física ocupada pela imagem do kernel.

A função \texttt{map\_page} em \texttt{kernel.c} gerencia a associação entre endereços virtuais e físicos em blocos alinhados a 4~KB:
\begin{verbatim}
void map_page(uint32_t *table1, uint32_t vaddr, paddr_t paddr, uint32_t flags) {
  if (!is_aligned(vaddr, PAGE_SIZE))
    PANIC("unaligned vaddr %x", vaddr);

  if (!is_aligned(paddr, PAGE_SIZE))
    PANIC("unaligned paddr %x", paddr);

  uint32_t vpn1 = (vaddr >> 22) & 0x3ff;
  if ((table1[vpn1] & PAGE_V) == 0) {
    uint32_t pt_paddr = alloc_pages(1);
    table1[vpn1] = ((pt_paddr / PAGE_SIZE) << 10) | PAGE_V;
  }

  uint32_t vpn0 = (vaddr >> 12) & 0x3ff;
  uint32_t *table0 = (uint32_t *)((table1[vpn1] >> 10) * PAGE_SIZE);
  table0[vpn0] = ((paddr / PAGE_SIZE) << 10) | flags | PAGE_V;
}
\end{verbatim}

Quando uma entrada de primeiro nível ainda não existe (\texttt{PAGE\_V == 0}), a função aloca dinamicamente uma página para a tabela de segundo nível através de \texttt{alloc\_pages(1)}. Em seguida, a entrada folha correspondente é gravada com o número da página física ($PPN = paddr / PAGE\_SIZE$), as flags de permissão e o bit de validade.

Na função \texttt{create\_process}, cada tarefa recebe uma tabela de páginas raiz. Para garantir que o código e as estruturas do kernel permaneçam acessíveis após a ativação da MMU, realiza-se o mapeamento identidade ($vaddr = paddr$) em toda a faixa de \texttt{\_\_kernel\_base} até \texttt{\_\_free\_ram\_end} com permissões completas (\texttt{PAGE\_R | PAGE\_W | PAGE\_X}):
\begin{verbatim}
  uint32_t *page_table = (uint32_t *)alloc_pages(1);
  for (paddr_t paddr = (paddr_t)__kernel_base;
       paddr < (paddr_t)__free_ram_end; paddr += PAGE_SIZE)
    map_page(page_table, paddr, paddr, PAGE_R | PAGE_W | PAGE_X);

  proc->page_table = page_table;
\end{verbatim}

Durante a troca de processos em \texttt{yield}, a MMU é atualizada gravando-se no registrador \texttt{satp} o identificador de modo \texttt{SATP\_SV32} combinado com o PPN da tabela raiz do próximo processo. A escrita no \texttt{satp} é envolvida por instruções \texttt{sfence.vma} antes e depois, garantindo a conclusão das alterações em memória e a invalidação de entradas obsoletas na TLB:
\begin{verbatim}
  __asm__ __volatile__(
      "sfence.vma\n"
      "csrw satp, %[satp]\n"
      "sfence.vma\n"
      "csrw sscratch, %[sscratch]\n"
      :
      : [satp] "r"(SATP_SV32 | ((uint32_t)next->page_table / PAGE_SIZE)),
        [sscratch] "r"((uint32_t)&next->stack[sizeof(next->stack)]));
\end{verbatim}

\section{Atividades Não Realizadas}

Mantendo a precisão sobre o escopo entregue, relata-se as tarefas pendentes do segundo bloco:

\begin{itemize}
    \item \textbf{Trabalhos 09 e 10 (Modo Usuário e Syscalls):} A segregação de privilégios com o nível U-Mode e a interface de chamadas de sistema (\texttt{ecall}) para viabilizar a execução segura de aplicações de usuário permanecem pendentes para as próximas etapas.
\end{itemize}

\section{Resultados e Validação}

O código foi compilado através do Clang/LLD com as opções \texttt{-ffreestanding -nostdlib -O2} e verificado no QEMU (\texttt{run.sh}).

\begin{enumerate}
    \item \textbf{Validação da Alocação de Páginas:} A inspeção dos símbolos via \texttt{llvm-nm} confirmou que \texttt{\_\_free\_ram} foi alocado em \texttt{0x80231000} e \texttt{\_\_free\_ram\_end} em \texttt{0x84231000}, respeitando o alinhamento de 4~KB e a extensão de 64~MB. A chamada de teste em \texttt{kernel\_main} gerou a saída serial esperada:
\begin{verbatim}
alloc_pages test: paddr0=80231000
alloc_pages test: paddr1=80233000
\end{verbatim}
    O endereço \texttt{paddr0} coincide com a base da memória livre, e \texttt{paddr1} situa-se exatamente 8192 bytes (2 páginas) à frente, confirmando a ausência de sobreposição.

    \item \textbf{Validação da Troca de Contexto e Escalonamento:} Foram criados dois processos concorrentes (\texttt{proc\_a} e \texttt{proc\_b}) que emitem repetidamente os caracteres \texttt{'A'} e \texttt{'B'}, cedendo a CPU voluntariamente via \texttt{yield}:
\begin{verbatim}
starting process A
Astarting process B
BABABABABABABABABABABABABABABABABABABABABABABABABABABABA...
\end{verbatim}
    A saída confirma que a alternância de fluxo ocorre de forma contínua e determinística, preservando o estado interno dos registradores e os ponteiros de pilha individuais de cada tarefa.

    \item \textbf{Validação da Paginação Sv32 e MMU:} A ativação da MMU e a integridade das tabelas foram comprovadas diretamente pelo monitor do QEMU:
    \begin{itemize}
        \item O comando \texttt{info registers} confirmou o registrador \texttt{satp = 80080246}, demonstrando que o bit 31 (\texttt{SATP\_SV32}) está ativo e que o PPN raiz aponta para o endereço físico correspondente da tabela de primeiro nível alocada.
        \item O comando \texttt{info mem} confirmou o mapeamento identidade integral do kernel de \texttt{0x80200000} até \texttt{0x84231000} com atributos de permissão \texttt{rwx}.
        \item A execução contínua dos processos sob paginação ativa ocorreu sem geração de falhas espúrias. Ademais, eventuais acessos inválidos a regiões não mapeadas são interceptados pelo tratador de interrupções (\texttt{handle\_trap}) via \texttt{stvec}, emitindo um pânico controlado sem reiniciar o processador.
    \end{itemize}
\end{enumerate}

\section{Conclusão}

As atividades desenvolvidas nesta segunda etapa dotaram o kernel de recursos fundamentais de infraestrutura: alocação dinâmica de memória física por páginas, concorrência cooperativa de processos com pilhas dedicadas e gerenciamento de memória virtual paginada sob a arquitetura Sv32 do RISC-V. A continuidade das tarefas envolverá o isolamento de aplicações no nível de privilégio U-Mode e a implementação do canal de chamadas de sistema.

\begin{thebibliography}{9}
\bibitem{nuta2023}
Nuta, S. (2023). \textit{Operating System in 1,000 Lines}.

\bibitem{sbispec}
RISC-V SBI Task Group (2022). \textit{RISC-V Supervisor Binary Interface Specification}.

\bibitem{riscvpriv}
Waterman, A. and Asanovi\'{c}, K. (2021). \textit{The RISC-V Instruction Set Manual, Volume II: Privileged Architecture}. RISC-V International.
\end{thebibliography}

\end{document}
